% =====================================================================
% Gemeinsame Präambel für main.tex (Konzept) und fragenkatalog.tex
% Einbinden mit % =====================================================================
% Gemeinsame Präambel für main.tex (Konzept) und fragenkatalog.tex
% Einbinden mit % =====================================================================
% Gemeinsame Präambel für main.tex (Konzept) und fragenkatalog.tex
% Einbinden mit % =====================================================================
% Gemeinsame Präambel für main.tex (Konzept) und fragenkatalog.tex
% Einbinden mit \input{praeambel}
% =====================================================================
\usepackage[utf8]{inputenc}
\usepackage[T1]{fontenc}
\usepackage{lmodern}
\usepackage[ngerman]{babel}
\usepackage[margin=2.3cm]{geometry}
\usepackage{enumitem}
\usepackage{booktabs}
\usepackage{tabularx}
\usepackage{xcolor}
\usepackage{listings}
\usepackage{fancyhdr}
\usepackage{amsmath}
\usepackage{array}
\usepackage[hidelinks]{hyperref}

% ---------- Farben ----------
\definecolor{zwetschge}{HTML}{5B2A86}
\definecolor{offen}{HTML}{C0392B}

% ---------- Kopf- und Fusszeile ----------
\pagestyle{fancy}
\fancyhf{}
\lhead{Zwätschgenabend -- Pubquiz}
\rhead{HSLU Elektrotechnik}
\cfoot{\thepage}
\setlength{\headheight}{14pt}

% ---------- Listen ----------
\setlist[itemize]{itemsep=0.2em, topsep=0.3em}

% ---------- Hilfsmakros ----------
% Markiert Punkte, die das OK noch klären muss
\newcommand{\offen}[1]{\textcolor{offen}{\textbf{[offen:} #1\textbf{]}}}

\lstset{
  basicstyle=\ttfamily\small,
  frame=single,
  rulecolor=\color{zwetschge!50},
  language=C,
  columns=fullflexible,
  xleftmargin=1em
}

% =====================================================================
\usepackage[utf8]{inputenc}
\usepackage[T1]{fontenc}
\usepackage{lmodern}
\usepackage[ngerman]{babel}
\usepackage[margin=2.3cm]{geometry}
\usepackage{enumitem}
\usepackage{booktabs}
\usepackage{tabularx}
\usepackage{xcolor}
\usepackage{listings}
\usepackage{fancyhdr}
\usepackage{amsmath}
\usepackage{array}
\usepackage[hidelinks]{hyperref}

% ---------- Farben ----------
\definecolor{zwetschge}{HTML}{5B2A86}
\definecolor{offen}{HTML}{C0392B}

% ---------- Kopf- und Fusszeile ----------
\pagestyle{fancy}
\fancyhf{}
\lhead{Zwätschgenabend -- Pubquiz}
\rhead{HSLU Elektrotechnik}
\cfoot{\thepage}
\setlength{\headheight}{14pt}

% ---------- Listen ----------
\setlist[itemize]{itemsep=0.2em, topsep=0.3em}

% ---------- Hilfsmakros ----------
% Markiert Punkte, die das OK noch klären muss
\newcommand{\offen}[1]{\textcolor{offen}{\textbf{[offen:} #1\textbf{]}}}

\lstset{
  basicstyle=\ttfamily\small,
  frame=single,
  rulecolor=\color{zwetschge!50},
  language=C,
  columns=fullflexible,
  xleftmargin=1em
}

% =====================================================================
\usepackage[utf8]{inputenc}
\usepackage[T1]{fontenc}
\usepackage{lmodern}
\usepackage[ngerman]{babel}
\usepackage[margin=2.3cm]{geometry}
\usepackage{enumitem}
\usepackage{booktabs}
\usepackage{tabularx}
\usepackage{xcolor}
\usepackage{listings}
\usepackage{fancyhdr}
\usepackage{amsmath}
\usepackage{array}
\usepackage[hidelinks]{hyperref}

% ---------- Farben ----------
\definecolor{zwetschge}{HTML}{5B2A86}
\definecolor{offen}{HTML}{C0392B}

% ---------- Kopf- und Fusszeile ----------
\pagestyle{fancy}
\fancyhf{}
\lhead{Zwätschgenabend -- Pubquiz}
\rhead{HSLU Elektrotechnik}
\cfoot{\thepage}
\setlength{\headheight}{14pt}

% ---------- Listen ----------
\setlist[itemize]{itemsep=0.2em, topsep=0.3em}

% ---------- Hilfsmakros ----------
% Markiert Punkte, die das OK noch klären muss
\newcommand{\offen}[1]{\textcolor{offen}{\textbf{[offen:} #1\textbf{]}}}

\lstset{
  basicstyle=\ttfamily\small,
  frame=single,
  rulecolor=\color{zwetschge!50},
  language=C,
  columns=fullflexible,
  xleftmargin=1em
}

% =====================================================================
\usepackage[utf8]{inputenc}
\usepackage[T1]{fontenc}
\usepackage{lmodern}
\usepackage[ngerman]{babel}
\usepackage[margin=2.3cm]{geometry}
\usepackage{enumitem}
\usepackage{booktabs}
\usepackage{tabularx}
\usepackage{xcolor}
\usepackage{listings}
\usepackage{fancyhdr}
\usepackage{amsmath}
\usepackage{array}
\usepackage[hidelinks]{hyperref}

% ---------- Farben ----------
\definecolor{zwetschge}{HTML}{5B2A86}
\definecolor{offen}{HTML}{C0392B}

% ---------- Kopf- und Fusszeile ----------
\pagestyle{fancy}
\fancyhf{}
\lhead{Zwätschgenabend -- Pubquiz}
\rhead{HSLU Elektrotechnik}
\cfoot{\thepage}
\setlength{\headheight}{14pt}

% ---------- Listen ----------
\setlist[itemize]{itemsep=0.2em, topsep=0.3em}

% ---------- Hilfsmakros ----------
% Markiert Punkte, die das OK noch klären muss
\newcommand{\offen}[1]{\textcolor{offen}{\textbf{[offen:} #1\textbf{]}}}

\lstset{
  basicstyle=\ttfamily\small,
  frame=single,
  rulecolor=\color{zwetschge!50},
  language=C,
  columns=fullflexible,
  xleftmargin=1em
}
